% AUTHOR: Amlal El Mahrouss
% PURPOSE: WG05.TN001: The Mathematics of Execution.

\documentclass[11pt,a4paper]{article}

\usepackage[utf8]{inputenc}
\usepackage[T1]{fontenc}
\usepackage{amsmath,amssymb,amsthm}
\usepackage{hyperref}
\usepackage[margin=1in]{geometry}

\title{The Mathematics of Execution.}
\author{Amlal El Mahrouss\\
\texttt{amlal@nekernel.org}}
\date{January 2026}

\begin{document}

\bf \maketitle

\begin{center}
	\rule[0.01cm]{17cm}{0.01cm}
\end{center}

\begin{center}
\begin{abstract}
This paper covers the algebraic properties and concepts of `The Execution Semantics' paper by El Mahrouss, A (2026). It is intended as an additional resource over the aforementioned paper.
We will define their concepts alongside their properties. We assume the reader has read the previous paper on 'The Execution Semantics' by El Mahrouss, A.
\end{abstract}
\end{center}

\begin{center}
	\rule[1cm]{17cm}{0.01cm}
\end{center}

\section{The Execution Product Formula}

Let $Pe(x, ..., z)$ be an execution product of variable arguments $x$, $...$, to $z$ defined in compatible and composable execution context $\mathbb{E}$. \\
Let the index $i$ denote the current execution domain of an execution product.
\\ \\ Consider the following formula:
\begin{equation}
Pe(x, ..., z) = \prod_{i=1}^{n}(x_{i} \cdot ... \cdot z_{i})+\mathbb{U}
\end{equation}
In which we define as the execution product, where $\mathbb{U}$ is the Unknown Execution of $Pe(x, ..., z)$, now defined as $g(\mathbb{E})$.

\subsection{Properties of $Pe$}

\begin{enumerate}

    \item $Pe$ as defined previously as the `execution product' shall always be valid within the execution context $\mathbb{E}$.
    \item The execution context $\mathbb{E}$ shall not denote $\varnothing$.

\end{enumerate}

\section{The Unknown Execution Property $\mathbb{U}$}

Let $\mathbb{U} \in Pe$ be the `Unknown Execution' as $\mathbb{U} = g(\mathbb{E})$.\\
Let $\mathbb{V} = \varnothing$ denote the `Empty Execution', with $\mathbb{V} \notin Pe$.

\subsection{Properties of $\mathbb{U}$}

\begin{enumerate}

    \item $\mathbb{U}$ shall not be equal to $\mathbb{V}$.
    \item The Execution Domain of $Pe$ shall not be equal to $\mathbb{V}$.

\end{enumerate}

\section{Conclusion}

Such properties are essential to define Execution Theory as a mathematical construct, as presented in El Mahrouss, A. (2026).

\section{References}

\begin{enumerate}
    \item El Mahrouss, A. (2026). The Execution Semantics: On Axioms, Domains, and Authority. Zenodo. https://doi.org/10.5281/zenodo.18470651
\end{enumerate}

\end{document}
